\chapter{Networking}

CharlotteOS deliberately departs from the traditional UNIX networking model.
Rather than treating TCP/IP as the foundation of all communication, CharlotteOS
treats networking as \textbf{an extension of the operating system's message passing
and capability model.}

The guiding principle is:

\begin{quote}
\textbf{Remote IPC is local IPC with a different transport.}
\end{quote}

\section{Why not sockets?}

The Berkeley socket API (\texttt{socket()}, \texttt{connect()}, \texttt{send()},
\texttt{recv()}) has served the Internet for decades. But it is not the right
foundation for a capability-based microkernel:

\begin{itemize}
    \item \textbf{Sockets identify hosts, not services.}
        \texttt{10.2.3.17:443} tells you where to connect, not what you are
        talking to. The mapping from address to service is out-of-band (DNS).

    \item \textbf{Sockets are byte streams, not messages.} Message boundaries
        are lost. Every protocol invents its own framing on top of TCP's
        undifferentiated stream.

    \item \textbf{Sockets carry ambient authority.} Any process can open a
        socket to any address. There is no capability check; authority comes
        from being able to reach the network, not from being granted a
        connection.

    \item \textbf{Sockets couple applications to transports.} Code written
        against TCP cannot transparently use shared memory for local
        communication or RDMA for high-performance communication.
\end{itemize}

CharlotteOS replaces all of this with a layered model where each layer has a
single responsibility.

The stack below is the target architecture. The current repository contains a
userspace virtio-net driver, a frame demultiplexer (``frouter''), a
reliable-message service with fragmentation, a distributed name service that
replicates a \texttt{name $\rightarrow$ \{node, generation\}} catalog across two guests, a
smoltcp TCP/IP service, and a remote-invocation path through the catalog. The
NIC/TCP path is runtime-validated on macOS QEMU TCG (\texttt{--relmsg-test},
\texttt{--disco-test}, \texttt{--dns-test}, \texttt{--tcpip-test}), and remote
scalar remote-call prototype is implemented (\texttt{dns::OP\_CALL}); general
remote capability delegation is not.

\section{The native protocol stack}

\begin{verbatim}
Applications
        |
Capability Invocation
        |
Distributed Objects
        |
RPC
        |
Reliable Message Layer
        |
Ethernet Frame Transport
        |
VirtIO / NIC Driver
\end{verbatim}

\subsection{Ethernet --- generic frame transport}

The NIC driver exports raw Ethernet frames. It knows about MAC addresses, MTU,
and frame CRCs. It knows \textbf{nothing} about IP, TCP, RPC, or services. Its
endpoint protocol has two operations:

\begin{itemize}
    \item \texttt{OP\_SEND(frame\_buffer)} --- transmit a frame via a moved
        memory object.
    \item \texttt{OP\_RECV} --- deferred receive: the caller retains a reply
        token; the driver replies with a received frame buffer.
\end{itemize}

Ethernet is a generic frame transport that can carry IP, ARP, LLDP, custom protocols, 
or CharlotteOS native messages. The driver provides the frames and higher layers 
interpret them.

\subsection{Reliable message layer}

This layer provides sequenced, acknowledged, reliable delivery of \textbf{messages}
(not streams) over the raw frame transport. It handles:

\begin{itemize}
    \item Sequencing (each message has a monotonic sequence number).
    \item Acknowledgements (the receiver confirms delivery).
    \item Retransmission (unacknowledged messages are resent).
    \item Fragmentation (messages larger than MTU are split and reassembled).
    \item Batching (multiple messages in one frame, one ACK for multiple messages).
    \item Congestion control (bounded in-flight windows).
\end{itemize}

The abstraction exported is a \textbf{reliable message}, not a byte stream.
This is the building block for RPC, consensus protocols, and distributed objects.

Fragmentation is implemented: a message larger than one frame's payload
($\approx$1460 bytes in wire protocol v2) is split across frames that share the message's sequence
number, each fragment carrying its byte offset and a \texttt{FLAG\_FRAG}/
\texttt{FLAG\_MORE} pair in the header's reserved bits. The receiver buffers
fragments by offset and reassembles the message once the last fragment
arrives. The application-level ceiling is \texttt{relmsg::MAX\_MSG}
(65535 bytes).

Wire protocol v2 carries a 64-bit sender-session identifier.  Data and
acknowledgement frames assert it with \texttt{FLAG\_SYN}.  Observing a new,
non-retired session resets the per-peer sequence spaces (including an
outstanding transmission), so one reliable-message service can restart while
its peer remains running.  A bounded retired-session window rejects delayed
frames from prior instances.  The receive queue and fragment count/bytes are
also bounded; a full delivery queue withholds sequence advancement and its
acknowledgement, applying retransmission-based backpressure.

\subsection{RPC layer}

The RPC layer provides request/reply semantics over the reliable message layer:

\begin{itemize}
    \item \textbf{Synchronous-looking futures:} \texttt{service.call(opcode, args).await?}
    \item \textbf{Asynchronous invocation:} the \texttt{Future} is backed by a
        pending RPC call; the shard continues other work while waiting.
    \item \textbf{Object references:} capabilities can be serialized and sent
        over RPC (delegation across machines).
    \item \textbf{Serialization:} typed protocol crates handle encoding/decoding
        between Rust types and the wire format.
\end{itemize}

The implemented DNS v2 remote-call path is narrower than this target model. It
currently carries scalar opcode/argument calls. A request is identified by
caller node, caller DNS-session generation, and call ID, and a reply is
accepted only from the expected peer. The implementation bounds outstanding
calls at 64, returns \texttt{ERR\_UNCERTAIN} after five seconds (execution may
already have occurred), and retains 128 completed results to suppress duplicate
execution within that window. The replicated catalog assigns a monotonically
advancing service generation; requests and replies carry it, and the target
rejects stale generations before execution. Remote object-capability
delegation remains future work.

\subsection{Distributed objects}

At the top of the stack, applications interact with \textbf{capabilities} ---
the same abstraction they use for local IPC:

\begin{verbatim}
let payroll: Capability<PayrollService> = lookup("Payroll")?;
payroll.call("calculate_salary", employee).await?;
\end{verbatim}

Whether the Payroll service runs on the same machine (local endpoint IPC) or on
a different machine (remote RPC over reliable messages over Ethernet) is
\textbf{transparent to the application}. The capability lookup returns a
connection; the transport layer behind that connection is an implementation
detail.

\section{TCP/IP as a compatibility service}

TCP/IP is supported as a userspace compatibility service, not the native
programming model. The frouter owns the NIC's deferred \texttt{OP\_RECV} and
routes IPv4 (\texttt{0x0800}) and ARP (\texttt{0x0806}) frames to the
\texttt{tcpip} service's \texttt{OP\_FRAME} ingress; a smoltcp
\texttt{phy::Device} adapter pulls those frames and transmits via
\texttt{OP\_SEND}:

\begin{verbatim}
// smoltcp adapter (receive is fed by the frouter's OP_FRAME ingress):
fn receive(&mut self, now: Instant) -> Option<(IpPacket, ...)> {
    self.rx.pop_front() // frames queued by the tcpip service
}

fn transmit(&mut self, now: Instant) -> Option<...> {
    // Calls OP_SEND with a memory object.
}
\end{verbatim}

The TCP/IP stack itself uses smoltcp 0.13. It is wrapped in a userspace
service that registers as \texttt{"tcpip"} and serves clients through a
socket-API protocol (\texttt{OP\_SOCKET}, \texttt{OP\_CONNECT},
\texttt{OP\_BIND}/\texttt{OP\_LISTEN}, \texttt{OP\_ACCEPT}, \texttt{OP\_SEND},
\texttt{OP\_RECV}, \texttt{OP\_CLOSE}). Legacy software that requires TCP/IP
connects to this service.

\section{Transport independence}

The reliable-message interface is designed to be transport-agnostic. Candidate
transports include:

\begin{itemize}
    \item \textbf{Ethernet} --- for inter-machine communication.
    \item \textbf{Shared memory} --- for co-located processes on the same machine,
        where the kernel can set up shared memory regions instead of going through
        the NIC.
    \item \textbf{Loopback} --- for same-machine same-stack communication.
    \item \textbf{RDMA} --- for high-performance cluster interconnects.
    \item \textbf{PCIe NTB} --- for backplane communication in multi-socket
        systems.
\end{itemize}

The intended RPC layer and everything above it should remain unchanged across
transports.
This is the same principle that makes local IPC transparent to the application
when it moves from same-process to same-machine to different-machine.

\section{Zero-copy networking}

The intended network stack avoids copying whenever practical:

\begin{itemize}
    \item \textbf{Receive path:} the NIC driver receives a buffer from hardware.
        That buffer (a memory object) is passed up the stack by transferring
        ownership --- first to the message layer (for reassembly), then to the
        RPC layer (for deserialization), then to the application. No intermediate
        copies.
    \item \textbf{Transmit path:} the application creates a buffer. It is passed
        down by moving ownership. The NIC driver DMA's it directly from the
        application's pages.
    \item \textbf{Buffer recycling:} after the application consumes a received
        buffer, it returns the (now empty) memory object to the driver's RX pool.
        The driver reuses it for DMA and the cycle repeats. IOMMU-backed pinning
        is not yet implemented.
\end{itemize}

This is the memory-object transfer model (Chapter~6) applied to networking.

\section{Why the networking architecture matters for critical software}

\begin{enumerate}
    \item \textbf{Remote IPC = local IPC.} An application that invokes a
        capability does not know or care whether the service is local or
        remote. This means:
        \begin{itemize}
            \item Services can be relocated without application changes.
            \item Testing can substitute a local mock for a remote service.
            \item A common interface can reduce environment-specific application
                code, although the network boundary still introduces failures
                absent from local IPC.
        \end{itemize}

    \item \textbf{Capability-based, not address-based.} An application invokes
        \texttt{PayrollService}, not \texttt{10.2.3.17:443}. Authority is a
        capability, not reachability. A compromised application that can
        reach the network cannot fabricate a connection to a service it was
        not granted.

    \item \textbf{Copy avoidance is a design goal.} Page ownership transfer can
        avoid intermediate payload copies. CPU mappings are MMU-enforced; DMA
        isolation awaits IOMMU/SMMU support. No performance comparison with
        sockets has yet been established.

    \item \textbf{Message boundaries are preserved.} The native abstraction is a
        message with identity, ownership, metadata, and boundaries. There is no
        need for an additional stream-framing layer. This can simplify protocol
        implementations, but does not eliminate parsing or framing defects.

    \item \textbf{TCP/IP is available but not mandatory.} Legacy software
        connects to the TCP/IP service. Native software uses the capability-based
        stack. The transition is incremental: a service can expose both a native
        endpoint and a TCP/IP socket through the compatibility service.
\end{enumerate}

\endinput
